\section{A longitudinal case: one covariance matrix, many inverse graphs}
\sectionmark{Many inverse graphs}

An essay written in 2022 under the deliberately exaggerated title ``Inverse covariance is all
you need'' placed a Gaussian chain, MSE reconstruction, sparse and ridge precision estimation,
network deconvolution, balanced network deconvolution, and eigenmode deletion in one sequence
\cite{wang2022inverse}.  The visual argument was compelling because the methods repeatedly took
the same form: diagonalize a matrix, alter its eigenvalues, and reconstruct a cleaner graph.  Read
through the operator framework, this family resemblance is real and stronger than a metaphor.
For a symmetric observed matrix
\begin{equation}
O=U\diag(\lambda_1,\ldots,\lambda_d)U^{\mathsf T},
\end{equation}
most of the methods in the essay can be written
\begin{equation}
O\longmapsto
U\diag\!\left(f(\lambda_1),\ldots,f(\lambda_d)\right)U^{\mathsf T},
\label{eq:inverse-covariance-spectral-family}
\end{equation}
possibly preceded by centering or scaling and followed by block compression and contact ranking.
The eigenvectors are often retained; the scientific decision is concentrated in the scalar
function $f$.

That algebraic statement explains why the works look so similar.  Each $f$ nevertheless estimates
the graph defined by a different forward story.  Covariance
inversion assumes a Gaussian probability law.  Network deconvolution assumes that indirect paths
are matrix powers of a direct network.  Balanced deconvolution changes which path lengths exist in
that story.  Leading-mode deletion treats selected spectral directions as nuisance rather than as
indirect paths.  A deep contact predictor may use the transformed matrix only as an input feature,
so its final edge is produced by the subsequent network rather than by $f$ alone.

The common family can therefore be typed by the chain
\begin{equation}
\boxed{
\begin{gathered}
\text{observed statistic}
\longrightarrow\text{propagation law}
\longrightarrow\text{inverse or spectral transform}\\
\longrightarrow\text{regularization}
\longrightarrow\text{edge object}
\longrightarrow\text{validation endpoint}
\end{gathered}}
\label{eq:inverse-graph-audit-chain}
\end{equation}
Two methods can be nearly identical in numerical linear algebra and different at any later arrow.
This section rewrites the historical sequence in those coordinates.

\subsection{The Gaussian chain: when the inverse really is the graph}

Let independent unit Gaussian noises generate
\begin{equation}
A=\epsilon_A,\qquad B=A+\epsilon_B,\qquad
C=B+\epsilon_C,\qquad D=C+\epsilon_D.
\label{eq:gaussian-four-chain}
\end{equation}
The covariance matrix is dense because an upstream fluctuation propagates to every downstream
variable.  For example, $A$ covaries with $C$ and $D$ even though the structural equations contain
no $A$--$C$ or $A$--$D$ term.  The joint negative log density, however, is
\begin{equation}
-\log p(a,b,c,d)=\frac12\left[a^2+(b-a)^2+(c-b)^2+(d-c)^2\right]+\text{constant}.
\end{equation}
Its Hessian, and hence its precision matrix $\Theta=\Sigma^{-1}$, is
\begin{equation}
\Theta=
\begin{pmatrix}
2&-1&0&0\\
-1&2&-1&0\\
0&-1&2&-1\\
0&0&-1&1
\end{pmatrix}.
\label{eq:gaussian-four-chain-precision}
\end{equation}
The tridiagonal support recovers the conditional-independence graph: for a nonsingular Gaussian
law, $\Theta_{ij}=0$ if and only if $X_i$ and $X_j$ are conditionally independent given all other
coordinates \cite{dempster1972}.  Here ``direct'' has an exact meaning because the probability
family, conditioning set, and node variables have all been declared.

This example also gives precision a limited physical interpretation.  In a Gaussian equilibrium
field, covariance measures fluctuations or susceptibility, whereas precision is the quadratic
restoring operator: it is the Hessian of the dimensionless Gaussian free energy.  A large
eigenvalue of $\Sigma$ is a soft, strongly fluctuating direction; inversion makes it a small
stiffness.  A small covariance eigenvalue is a stiff direction and becomes a large precision
eigenvalue.  This language becomes literal molecular physics only if the variables are physical
coordinates, the distribution is an equilibrium ensemble, and hidden degrees of freedom have
been handled.  For one-hot MSA columns, $\Theta$ is instead a Gaussian surrogate for categorical
sequence statistics.  Its graph is probability-exact inside that surrogate; forces and contacts
require an additional correspondence.

\subsection{MSE reconstruction: the same optimum through another objective}

Chapter~7 showed that a regularized linear self-reconstruction problem can have an optimum that is
an affine function of $(S+\lambda I)^{-1}$.  The 2022 essay emphasized the same fact from the
opposite direction: a matrix learned by squared-error reconstruction can land on the same
precision-like solution as a multivariate Gaussian fit even though the loss surfaces are not the
same.  One route differentiates a log determinant and imposes positive definiteness; the other
differentiates a quadratic reconstruction loss and can be implemented as an ordinary linear
layer.

The exact shared optimum is valuable.  It permits solvers and regularizers to move between
statistical modeling and neural reconstruction.  The Gaussian objective identifies a normalized
joint density when its constraints are satisfied.
The MSE objective identifies a best linear reconstruction map under a chosen data geometry.  A
common stationary point makes the fitted matrices equal, while thermodynamic and probabilistic
interpretations remain tied to their respective objectives and constraints.

This is the same pattern encountered in the 2019 unified framework \cite{dauparas2019}.  The
horizontal equivalence there grouped models by reconstruction loss.  The present longitudinal
case groups them by spectral action.  Together they show why implementation-level unification is
scientifically productive: it exposes exactly which remainder must still be typed.

\subsection{Sparse and ridge precision: regularization chooses the graph}

Given an empirical covariance $S$, regularized Gaussian precision estimation takes the generic
form
\begin{equation}
\widehat\Theta
=\arg\min_{\Theta\succ0}
\left\{\operatorname{tr}(S\Theta)-\log\det\Theta+R(\Theta)\right\}.
\label{eq:regularized-gaussian-precision}
\end{equation}
PSICOV used sparse inverse covariance estimation for protein contact prediction
\cite{jones2012}.  An $L_1$ penalty makes many fitted entries exactly zero, so it performs two jobs:
it stabilizes an underdetermined estimate and declares a sparse Gaussian conditional graph.  Its
zeros belong to that regularized model and express its selected conditional graph.  Experimental
influence remains a separate endpoint.

ResPRE instead used the Frobenius penalty $R(\Theta)=\rho\|\Theta\|_F^2$ and then supplied the raw
$21\times21$ precision blocks to a deep residual contact predictor \cite{li2019respre}.  Its
stationarity equation is
\begin{equation}
S-\Theta^{-1}+2\rho\Theta=0.
\label{eq:respre-stationarity}
\end{equation}
If $S=V\Lambda V^{\mathsf T}$, the solution shares the eigenvectors of $S$ and transforms each
eigenvalue by
\begin{equation}
k_\rho(\lambda)
=\frac{-\lambda+\sqrt{\lambda^2+8\rho}}{4\rho},
\qquad
\widehat\Theta=V\diag(k_\rho(\lambda_i))V^{\mathsf T}.
\label{eq:respre-ridge-filter}
\end{equation}
Thus the $L_2$ estimator is literally a smooth spectral filter.  It usually produces graded
shrinkage, while $L_1$ fitting can select an exact support.  ResPRE then continued beyond the
precision graph.
Its final edge was a contact probability learned by a residual convolutional network from the
full categorical blocks.  The precision matrix was a representation that made indirect
correlations easier for a finite neural predictor to use.  Indeed, covariance and invertible
precision contain the same information in principle; their empirical difference arose because
the downstream learner, sample size, and regularization were finite.

The two papers therefore differ less in their observed statistic than in the composite estimator.
PSICOV couples Gaussian fitting to sparse support and a contact score.  ResPRE couples ridge
spectral shrinkage to a nonlinear, structurally supervised decoder.  Calling both ``inverse
covariance methods'' is correct but incomplete in the same way that calling AlphaFold 1 and
AlphaFold 2 ``distance predictors'' is correct but incomplete.

\subsection{Network deconvolution: inversion of a walk algebra}

Network deconvolution begins from a different forward model.  It assumes that an observed
similarity network contains direct edges and all of their transitive walks:
\begin{equation}
G_{\mathrm{obs}}
=G_{\mathrm{dir}}+G_{\mathrm{dir}}^2+G_{\mathrm{dir}}^3+\cdots
=G_{\mathrm{dir}}(I-G_{\mathrm{dir}})^{-1}.
\label{eq:network-deconvolution-forward}
\end{equation}
Under the convergence and diagonalizability assumptions, inversion gives
\begin{equation}
G_{\mathrm{dir}}
=G_{\mathrm{obs}}(I+G_{\mathrm{obs}})^{-1},
\qquad
f_{\mathrm{ND}}(\lambda)=\frac{\lambda}{1+\lambda}
\label{eq:network-deconvolution-filter}
\end{equation}
after the paper's required spectral scaling \cite{feizi2013}.  The expression resembles a
precision transform because both contain a matrix inverse.  The inverse is solving a different
equation.  Gaussian precision inverts a fluctuation covariance under a joint probability law;
network deconvolution inverts a geometric series of walks under a propagation model.  Its
off-diagonal entry means ``direct relative to the assumed walk closure,'' not ``conditionally
independent in a Gaussian law.''

The distinction is more than philosophical.  Matrix multiplication counts weighted walks only
after a node basis, edge composition rule, and treatment of signs have been fixed.  A categorical
operator graph requires compatible fiber maps before $G^2$ even has a meaning.  For an ordinary
scalar symmetric network the composition is automatic, but its biological validity remains an
empirical assumption.  Network deconvolution improved several biological networks because that
assumption was useful; the algebra alone cannot prove that an observed covariance was generated
by the full walk series.

\subsection{Balanced deconvolution: changing which paths exist}

Balanced network deconvolution changed the forward model by retaining only odd path lengths:
\begin{equation}
G_{\mathrm{obs}}
=G_{\mathrm{dir}}+G_{\mathrm{dir}}^3+G_{\mathrm{dir}}^5+\cdots
=G_{\mathrm{dir}}(I-G_{\mathrm{dir}}^2)^{-1}.
\label{eq:bnd-forward}
\end{equation}
Consequently its spectral inverse is
\begin{equation}
f_{\mathrm{BND}}(\lambda)
=\begin{cases}
\dfrac{-1+\sqrt{1+4\lambda^2}}{2\lambda},&\lambda\ne0,\\[6pt]
0,&\lambda=0.
\end{cases}
\label{eq:bnd-filter}
\end{equation}
Removing even powers preserves the sign balance of eigenvalues and maps every real observed
eigenvalue to a direct eigenvalue with magnitude below one, avoiding the network-dependent scaling
required by the original transform \cite{sun2015}.  Applied as a postprocessor, BND improved
contact rankings from several coevolutionary methods, with smaller gains for methods such as DCA
and PSICOV that had already attempted to suppress transitive effects.

In the language of this book, BND changes the algebra of admissible paths.  It should not be
described as the same inverse covariance estimator with a technical fix: even-length paths are
absent from its generative story.  Nor does odd-only propagation have an automatic physical law
behind it.  Its justification is a combination of spectral behavior and downstream validation.
This is a legitimate scientific route, but the validation endpoint, rather than the square root in
Eq.~\ref{eq:bnd-filter}, supplies the biological evidence.

\subsection{Leading-mode deletion: a nuisance model, not an inverse}

Qin and Colwell analyzed why phylogenetic relatedness contaminates covariance matrices built from
MSAs \cite{qin2018}.  Their random-matrix calculation predicts a power-law tail among the largest
covariance eigenvalues.  On that model, selected leading directions are dominated by shared
ancestry rather than by structural coupling.  The corresponding correction is a spectral
projection such as
\begin{equation}
C_{\mathrm{clean}}
=C-\sum_{k\in\mathcal P}\lambda_k u_ku_k^{\mathsf T},
\label{eq:colwell-leading-mode-deletion}
\end{equation}
where $\mathcal P$ is chosen from the declared phylogenetic-tail criterion.  Removing those modes
improved contact prediction in their tests.

This method is especially revealing because it acts on the same eigendecomposition as precision
estimation to target a different object from $C^{-1}$.  Inversion applies $f(\lambda)=1/\lambda$ to every
retained direction, suppressing large covariance modes while amplifying small ones.  Leading-mode
deletion applies $f(\lambda)=0$ to a selected nuisance set and $f(\lambda)=\lambda$ elsewhere.  The
first transform follows from a Gaussian conditional model; the second follows from a phylogenetic
null model.  They can improve the same contact benchmark for different reasons.

This corrects an overstatement in the 2022 essay.  The Colwell construction belongs in the same
spectral family.  Its conceptual advance is to identify which large modes should be removed and
supply a model for
why.  Without that null, deleting leading eigenvectors is only rank-relative cleaning and may
erase real subfamily or functional structure.  With it, the projection is an estimand-level
decision of the kind developed in Chapter~9.

\subsection{What is genuinely the same}

The longitudinal sequence can now be summarized without either exaggerating or dissolving the
unification:
\begin{center}
\scriptsize
\begin{tabular}{@{}>{\raggedright\arraybackslash}p{0.14\linewidth}
                >{\raggedright\arraybackslash}p{0.26\linewidth}
                >{\raggedright\arraybackslash}p{0.23\linewidth}
                >{\raggedright\arraybackslash}p{0.29\linewidth}@{}}
\toprule
Method & Forward claim & Spectral action & Output semantics and endpoint \\
\midrule
Gaussian precision & covariance of one Gaussian joint law & $1/\lambda$ & conditional graph in the fitted Gaussian family \\
MSE reconstruction & best linear reconstruction under squared loss & regularized inverse, up to an affine identity term & reconstruction weight; Gaussian meaning only under the equivalence conditions \\
PSICOV / ResPRE & high-dimensional Gaussian surrogate plus regularization & sparse support or ridge filter $k_\rho$ & sparse coupling score, or full precision feature decoded into contacts \\
Network deconvolution & observation is the sum of all directed walk powers & $\lambda/(1+\lambda)$ after scaling & direct edge relative to the all-walk propagation model \\
BND & observation is the sum of odd walk powers & $f_{\mathrm{BND}}(\lambda)$ & direct edge relative to an odd-walk model; validated as contact postprocessing \\
Qin--Colwell & leading covariance tail is enriched for phylogenetic nuisance & delete selected leading modes & phylogeny-corrected covariance used for contact ranking \\
\bottomrule
\end{tabular}
\end{center}

At the level of Eq.~\ref{eq:inverse-covariance-spectral-family}, the methods really are close.
They reuse the same matrix, eigenvectors, rational or thresholded scalar functions, and often the
same contact benchmark.  This is why one implementation can suggest another and why
regularization transfers are fruitful.  At the level of scientific objects, their differences are
concentrated rather than negligible.  The forward law determines what has been inverted; the null
determines what has been removed; the regularizer determines what can be identified from finite
data; and the downstream decoder determines whether the transformed matrix is the answer or only
an input feature.

Formally, let a method be the composite map
\begin{equation}
\mathcal A_m
=\mathcal V_m\circ\mathcal E_m\circ\mathcal R_m\circ
\Phi_m\circ R_m(O-O_{0,m})R_m^{\mathsf T},
\label{eq:composite-inverse-graph-method}
\end{equation}
where $O_{0,m}$ is a null expectation, $R_m$ a nuisance residualizer, $\Phi_m$ the inverse or
spectral transform, $\mathcal R_m$ regularization, $\mathcal E_m$ edge extraction, and
$\mathcal V_m$ validation or decoding.  Two methods are fully equivalent on a declared class of
inputs only when these composite maps agree on that class.  Sharing $O$ or even sharing
$\Phi_m(O)$ is a weaker correspondence.  Agreement of final contact rankings is weaker still,
because scalarization can hide different categorical operators.

The durable conclusion is therefore not that inverse covariance is literally everything.  It is
that a surprisingly large body of protein inference can be organized as competing inverse
problems on one second-order statistic.  Once written that way, what looked like many unrelated
algorithms becomes one design space.  The important questions are no longer which paper used an
inverse, a square root, or an eigenvalue cutoff.  They are: which propagation law was assumed,
which modes were declared nuisance, which edge object survived the transform, and which experiment
made that declaration credible.
